\subsubsection*{Arquivo: src/server.js}
\begin{lstlisting}[caption={src/server.js}]
// src/server.js
require('dotenv').config();
const express = require('express');
const cors = require('cors');
const path = require('path');
const pool = require('./config/db');
const errorHandler = require('./middlewares/error.middleware');
const asyncHandler = require('./middlewares/async.middleware');
const authRoutes = require('./modules/auth/auth.routes');
const catalogRoutes = require('./modules/catalog/catalog.routes');
const orderRoutes = require('./modules/order/order.routes');
const notificationRoutes = require('./modules/notification/notification.routes');
const reportRoutes = require('./modules/report/report.routes');

const app = express();

app.use(cors());
app.use(express.json());
app.use(express.static(path.join(__dirname, '../frontend')));

app.use((req, res, next) => {
    console.log(`${req.method} ${req.url}`);
    next();
});

app.get('/api/health', asyncHandler(async (req, res) => {
  await pool.query('SELECT 1');
  res.status(200).json({ status: 'UP', database: 'connected' });
}));


app.use('/api/auth', authRoutes);
app.use('/api', catalogRoutes);
app.use('/api', orderRoutes);
app.use('/api/notificacoes', notificationRoutes);
app.use('/api/relatorios', reportRoutes);

app.use(errorHandler);

const PORT = process.env.PORT || 3000;
app.listen(PORT, () => {
  console.log(`Server is running on port ${PORT}`);
});\end{lstlisting}
\vspace{0.5cm}
\subsubsection*{Arquivo: src/config/db.js}
\begin{lstlisting}[caption={src/config/db.js}]
require('dotenv').config();
const { Pool } = require('pg');

const pool = new Pool({
  host: process.env.DB_HOST,
  port: process.env.DB_PORT,
  database: process.env.DB_NAME,
  user: process.env.DB_USER,
  password: process.env.DB_PASS,
});

module.exports = pool;
\end{lstlisting}
\vspace{0.5cm}
\subsubsection*{Arquivo: src/config/mailer.js}
\begin{lstlisting}[caption={src/config/mailer.js}]
// src/config/mailer.js
const nodemailer = require('nodemailer');

const transporter = nodemailer.createTransport({
  host: process.env.SMTP_HOST,
  port: process.env.SMTP_PORT,
  auth: {
    user: process.env.SMTP_USER,
    pass: process.env.SMTP_PASS,
  },
});

const enviarEmailNotificacao = async ({ destinatario, assunto, texto, html }) => {
  if (!process.env.SMTP_HOST) {
    console.warn('[Mailer] SMTP não configurado. Disparo ignorado com segurança.');
    return false;
  }

  try {
    await transporter.sendMail({
      from: process.env.SMTP_USER || '"SISFEIRA" <no-reply@sisfeira.local>',
      to: destinatario,
      subject: assunto,
      text: texto,
      html: html,
    });
    return true;
  } catch (err) {
    console.error('[Mailer Error]:', err.message);
    return false;
  }
};

module.exports = { enviarEmailNotificacao };\end{lstlisting}
\vspace{0.5cm}
\subsubsection*{Arquivo: src/middlewares/async.middleware.js}
\begin{lstlisting}[caption={src/middlewares/async.middleware.js}]
const asyncHandler = (fn) => (req, res, next) => {
  return Promise.resolve(fn(req, res, next)).catch(next);
};

module.exports = asyncHandler;
\end{lstlisting}
\vspace{0.5cm}
\subsubsection*{Arquivo: src/middlewares/auth.middleware.js}
\begin{lstlisting}[caption={src/middlewares/auth.middleware.js}]
// src/middlewares/auth.middleware.js
const jwt = require('jsonwebtoken');

const verifyToken = (req, res, next) => {
  const authHeader = req.headers['authorization'];
  const token = authHeader && authHeader.split(' ')[1];

  if (!token) {
    return res.status(401).json({
      erro: true,
      mensagem: 'Token inválido ou não fornecido',
      codigo: 'UNAUTHORIZED'
    });
  }

  try {
    const decoded = jwt.verify(token, process.env.JWT_SECRET || 'supersecret');
    req.user = decoded;
    next();
  } catch (error) {
    return res.status(401).json({
      erro: true,
      mensagem: 'Token inválido ou não fornecido',
      codigo: 'UNAUTHORIZED'
    });
  }
};

const requireRole = (rolesPermitidas) => {
  return (req, res, next) => {
    if (!req.user || !rolesPermitidas.includes(req.user.perfil)) {
      return res.status(403).json({
        erro: true,
        mensagem: 'Acesso não autorizado para este perfil',
        codigo: 'FORBIDDEN'
      });
    }
    next();
  };
};

module.exports = { verifyToken, requireRole };\end{lstlisting}
\vspace{0.5cm}
\subsubsection*{Arquivo: src/middlewares/error.middleware.js}
\begin{lstlisting}[caption={src/middlewares/error.middleware.js}]
const errorHandler = (err, req, res, next) => {
  const statusCode = err.statusCode || 500;
  res.status(statusCode).json({
    erro: true,
    mensagem: err.message || 'Erro interno do servidor',
    codigo: err.code || 'INTERNAL_SERVER_ERROR'
  });
};

module.exports = errorHandler;
\end{lstlisting}
\vspace{0.5cm}
\subsubsection*{Arquivo: src/modules/auth/auth.controller.js}
\begin{lstlisting}[caption={src/modules/auth/auth.controller.js}]
// src/modules/auth/auth.controller.js
const authService = require('./auth.service');
const asyncHandler = require('../../middlewares/async.middleware');

class AuthController {
  cadastrar = asyncHandler(async (req, res) => {
    const novoUsuario = await authService.cadastrar(req.body);
    res.status(201).json(novoUsuario);
  });

  login = asyncHandler(async (req, res) => {
    const credenciais = await authService.login(req.body);
    res.status(200).json(credenciais);
  });
}

module.exports = new AuthController();\end{lstlisting}
\vspace{0.5cm}
\subsubsection*{Arquivo: src/modules/auth/auth.repository.js}
\begin{lstlisting}[caption={src/modules/auth/auth.repository.js}]
const pool = require('../../config/db');

class AuthRepository {
  async buscarPorEmail(email) {
    const query = `
      SELECT id, nome, email, senha_hash, telefone, perfil 
      FROM usuarios 
      WHERE email = $1
    `;
    const result = await pool.query(query, [email]);
    return result.rows[0];
  }

  async criarUsuario({ nome, email, senha_hash, telefone, perfil }) {
    const query = `
      INSERT INTO usuarios (nome, email, senha_hash, telefone, perfil) 
      VALUES ($1, $2, $3, $4, $5) 
      RETURNING id, nome, email, telefone, perfil
    `;
    const values = [nome, email, senha_hash, telefone, perfil];
    const result = await pool.query(query, values);
    return result.rows[0];
  }
}

module.exports = new AuthRepository();
\end{lstlisting}
\vspace{0.5cm}
\subsubsection*{Arquivo: src/modules/auth/auth.routes.js}
\begin{lstlisting}[caption={src/modules/auth/auth.routes.js}]
// src/modules/auth/auth.routes.js
const express = require('express');
const authController = require('./auth.controller');

const router = express.Router();

router.post('/cadastro', authController.cadastrar);
router.post('/login', authController.login);

module.exports = router;\end{lstlisting}
\vspace{0.5cm}
\subsubsection*{Arquivo: src/modules/auth/auth.service.js}
\begin{lstlisting}[caption={src/modules/auth/auth.service.js}]
const bcrypt = require('bcrypt');
const jwt = require('jsonwebtoken');
const authRepository = require('./auth.repository');

class AuthService {
  async cadastrar({ nome, email, senha, telefone, perfil }) {
    if (!nome || !email || !senha || !perfil) {
      const error = new Error('Campos obrigatórios: nome, email, senha e perfil');
      error.statusCode = 400;
      throw error;
    }

    const perfisValidos = ['CLIENTE', 'PRODUTOR', 'ORGANIZACAO'];
    if (!perfisValidos.includes(perfil)) {
      const error = new Error('Perfil inválido');
      error.statusCode = 400;
      throw error;
    }

    if (perfil === 'PRODUTOR' && !telefone) {
      const error = new Error('Telefone é obrigatório para produtores');
      error.statusCode = 400;
      throw error;
    }

    const usuarioExistente = await authRepository.buscarPorEmail(email);
    if (usuarioExistente) {
      const error = new Error('E-mail já cadastrado');
      error.statusCode = 409;
      throw error;
    }

    const salt = await bcrypt.genSalt(10);
    const senha_hash = await bcrypt.hash(senha, salt);

    const novoUsuario = await authRepository.criarUsuario({
      nome,
      email,
      senha_hash,
      telefone: telefone || null,
      perfil
    });

    return novoUsuario;
  }

  async login({ email, senha }) {
    if (!email || !senha) {
      const error = new Error('Email e senha são obrigatórios');
      error.statusCode = 400;
      throw error;
    }

    const usuario = await authRepository.buscarPorEmail(email);
    if (!usuario) {
      const error = new Error('Credenciais inválidas');
      error.statusCode = 401;
      throw error;
    }

    const isMatch = await bcrypt.compare(senha, usuario.senha_hash);
    if (!isMatch) {
      const error = new Error('Credenciais inválidas');
      error.statusCode = 401;
      throw error;
    }

    const payload = {
      id: usuario.id,
      nome: usuario.nome,
      perfil: usuario.perfil
    };

    const token = jwt.sign(
      payload, 
      process.env.JWT_SECRET || 'supersecret', 
      { expiresIn: process.env.JWT_EXPIRES_IN || '8h' }
    );

    return {
      token,
      usuario: payload
    };
  }
}

module.exports = new AuthService();
\end{lstlisting}
\vspace{0.5cm}
\subsubsection*{Arquivo: src/modules/catalog/catalog.controller.js}
\begin{lstlisting}[caption={src/modules/catalog/catalog.controller.js}]
// src/modules/catalog/catalog.controller.js
const catalogService = require('./catalog.service');
const asyncHandler = require('../../middlewares/async.middleware');

class CatalogController {
  listarCatalogo = asyncHandler(async (req, res) => {
    const produtos = await catalogService.listarCatalogo();
    res.status(200).json(produtos);
  });

  listarProdutosDoProdutor = asyncHandler(async (req, res) => {
    const produtos = await catalogService.listarProdutosDoProdutor(req.user.id);
    res.status(200).json(produtos);
  });

  criarProduto = asyncHandler(async (req, res) => {
    const produto = await catalogService.criarProduto(req.user.id, req.body);
    res.status(201).json(produto);
  });

  atualizarProduto = asyncHandler(async (req, res) => {
    const { id } = req.params;
    const produto = await catalogService.atualizarProduto(id, req.user.id, req.body);
    res.status(200).json(produto);
  });

  desativarProduto = asyncHandler(async (req, res) => {
    const { id } = req.params;
    const resultado = await catalogService.desativarProduto(id, req.user.id);
    res.status(200).json(resultado);
  });
}

module.exports = new CatalogController();\end{lstlisting}
\vspace{0.5cm}
\subsubsection*{Arquivo: src/modules/catalog/catalog.repository.js}
\begin{lstlisting}[caption={src/modules/catalog/catalog.repository.js}]
// src/modules/catalog/catalog.repository.js
const pool = require('../../config/db');

class CatalogRepository {
  async listarCatalogoPublico() {
    const query = `
      SELECT p.id, p.nome, p.descricao, p.preco, p.unidade_medida, 
             u.nome AS produtor_nome, u.id AS produtor_id
      FROM produtos p
      INNER JOIN usuarios u ON p.produtor_id = u.id
      WHERE p.ativo = TRUE
      ORDER BY p.nome ASC
    `;
    const result = await pool.query(query);
    return result.rows;
  }

  async listarPorProdutor(produtorId) {
    const query = `
      SELECT id, nome, descricao, preco, unidade_medida, ativo 
      FROM produtos 
      WHERE produtor_id = $1 AND ativo = TRUE 
      ORDER BY nome ASC
    `;
    const result = await pool.query(query, [produtorId]);
    return result.rows;
  }

  async buscarPorId(id) {
    const query = `SELECT * FROM produtos WHERE id = $1`;
    const result = await pool.query(query, [id]);
    return result.rows[0];
  }

  async criarProduto({ produtor_id, nome, descricao, preco, unidade_medida }) {
    const query = `
      INSERT INTO produtos (produtor_id, nome, descricao, preco, unidade_medida) 
      VALUES ($1, $2, $3, $4, $5) 
      RETURNING *
    `;
    const values = [produtor_id, nome, descricao, preco, unidade_medida];
    const result = await pool.query(query, values);
    return result.rows[0];
  }

  async atualizarProduto(id, produtorId, { nome, descricao, preco, unidade_medida }) {
    const query = `
      UPDATE produtos 
      SET nome = $1, descricao = $2, preco = $3, unidade_medida = $4 
      WHERE id = $5 AND produtor_id = $6 
      RETURNING *
    `;
    const values = [nome, descricao, preco, unidade_medida, id, produtorId];
    const result = await pool.query(query, values);
    return result.rows[0];
  }

  async desativarProduto(id, produtorId) {
    const query = `
      UPDATE produtos 
      SET ativo = FALSE 
      WHERE id = $1 AND produtor_id = $2 
      RETURNING id
    `;
    const result = await pool.query(query, [id, produtorId]);
    return result.rows[0];
  }
}

module.exports = new CatalogRepository();\end{lstlisting}
\vspace{0.5cm}
\subsubsection*{Arquivo: src/modules/catalog/catalog.routes.js}
\begin{lstlisting}[caption={src/modules/catalog/catalog.routes.js}]
// src/modules/catalog/catalog.routes.js
const express = require('express');
const catalogController = require('./catalog.controller');
const { verifyToken, requireRole } = require('../../middlewares/auth.middleware');

const router = express.Router();

router.get('/catalogo', catalogController.listarCatalogo);
router.get('/produtos/produtor', verifyToken, requireRole(['PRODUTOR']), catalogController.listarProdutosDoProdutor);
router.post('/produtos', verifyToken, requireRole(['PRODUTOR']), catalogController.criarProduto);
router.put('/produtos/:id', verifyToken, requireRole(['PRODUTOR']), catalogController.atualizarProduto);
router.delete('/produtos/:id', verifyToken, requireRole(['PRODUTOR']), catalogController.desativarProduto);

module.exports = router;\end{lstlisting}
\vspace{0.5cm}
\subsubsection*{Arquivo: src/modules/catalog/catalog.service.js}
\begin{lstlisting}[caption={src/modules/catalog/catalog.service.js}]
// src/modules/catalog/catalog.service.js
const catalogRepository = require('./catalog.repository');

class CatalogService {
  async listarCatalogo() {
    const produtos = await catalogRepository.listarCatalogoPublico();
    return produtos.map(p => ({
      ...p,
      preco: Number(p.preco)
    }));
  }

  async listarProdutosDoProdutor(produtorId) {
    const produtos = await catalogRepository.listarPorProdutor(produtorId);
    return produtos.map(p => ({
      ...p,
      preco: Number(p.preco)
    }));
  }

  async criarProduto(produtorId, dados) {
    const { nome, descricao, preco, unidade_medida } = dados;

    if (!nome || preco === undefined || !unidade_medida) {
      const error = new Error('Campos obrigatórios: nome, preco e unidade_medida');
      error.statusCode = 400;
      throw error;
    }

    if (Number(preco) < 0) {
      const error = new Error('O preço não pode ser negativo');
      error.statusCode = 400;
      throw error;
    }

    return await catalogRepository.criarProduto({
      produtor_id: produtorId,
      nome,
      descricao: descricao || null,
      preco,
      unidade_medida
    });
  }

  async atualizarProduto(id, produtorId, dados) {
    const produtoExistente = await catalogRepository.buscarPorId(id);
    
    if (!produtoExistente) {
      const error = new Error('Produto não encontrado');
      error.statusCode = 404;
      throw error;
    }

    if (produtoExistente.produtor_id !== produtorId) {
      const error = new Error('Acesso negado para modificar este produto');
      error.statusCode = 403;
      throw error;
    }

    const { nome, descricao, preco, unidade_medida } = dados;

    return await catalogRepository.atualizarProduto(id, produtorId, {
      nome: nome || produtoExistente.nome,
      descricao: descricao !== undefined ? descricao : produtoExistente.descricao,
      preco: preco !== undefined ? preco : produtoExistente.preco,
      unidade_medida: unidade_medida || produtoExistente.unidade_medida
    });
  }

  async desativarProduto(id, produtorId) {
    const produtoExistente = await catalogRepository.buscarPorId(id);
    
    if (!produtoExistente) {
      const error = new Error('Produto não encontrado');
      error.statusCode = 404;
      throw error;
    }

    if (produtoExistente.produtor_id !== produtorId) {
      const error = new Error('Acesso negado para desativar este produto');
      error.statusCode = 403;
      throw error;
    }

    await catalogRepository.desativarProduto(id, produtorId);
    return { mensagem: 'Produto desativado com sucesso' };
  }
}

module.exports = new CatalogService();\end{lstlisting}
\vspace{0.5cm}
\subsubsection*{Arquivo: src/modules/notification/notification.controller.js}
\begin{lstlisting}[caption={src/modules/notification/notification.controller.js}]
// src/modules/notification/notification.controller.js
const notificationService = require('./notification.service');
const asyncHandler = require('../../middlewares/async.middleware');

class NotificationController {
  listar = asyncHandler(async (req, res) => {
    const notificacoes = await notificationService.listarNotificacoes(req.user.id);
    res.status(200).json(notificacoes);
  });

  marcarComoLida = asyncHandler(async (req, res) => {
    const { id } = req.params;
    const produtorId = req.user.id;
    const resultado = await notificationService.marcarLida(id, produtorId);
    res.status(200).json(resultado);
  });
}

module.exports = new NotificationController();\end{lstlisting}
\vspace{0.5cm}
\subsubsection*{Arquivo: src/modules/notification/notification.repository.js}
\begin{lstlisting}[caption={src/modules/notification/notification.repository.js}]
// src/modules/notification/notification.repository.js
const pool = require('../../config/db');

class NotificationRepository {
  async listarPorProdutor(produtorId) {
    const query = `
      SELECT n.id, n.pedido_id, n.mensagem, n.lida,
             p.comprovante_codigo, p.criado_em AS data_pedido
      FROM notificacoes n
      JOIN pedidos p ON n.pedido_id = p.id
      WHERE n.produtor_id = $1
      ORDER BY n.lida ASC, p.criado_em DESC;
    `;
    const result = await pool.query(query, [produtorId]);
    return result.rows;
  }

  async marcarComoLida(notificacaoId, produtorId) {
    const query = `
      UPDATE notificacoes 
      SET lida = TRUE 
      WHERE id = $1 AND produtor_id = $2 
      RETURNING id, lida;
    `;
    const result = await pool.query(query, [notificacaoId, produtorId]);
    return result.rows[0];
  }
}

module.exports = new NotificationRepository();\end{lstlisting}
\vspace{0.5cm}
\subsubsection*{Arquivo: src/modules/notification/notification.routes.js}
\begin{lstlisting}[caption={src/modules/notification/notification.routes.js}]
// src/modules/notification/notification.routes.js
const express = require('express');
const notificationController = require('./notification.controller');
const { verifyToken, requireRole } = require('../../middlewares/auth.middleware');

const router = express.Router();

router.get('/produtor', verifyToken, requireRole(['PRODUTOR']), notificationController.listar);
router.patch('/:id/lida', verifyToken, requireRole(['PRODUTOR']), notificationController.marcarComoLida);

module.exports = router;\end{lstlisting}
\vspace{0.5cm}
\subsubsection*{Arquivo: src/modules/notification/notification.service.js}
\begin{lstlisting}[caption={src/modules/notification/notification.service.js}]
// src/modules/notification/notification.service.js
const notificationRepository = require('./notification.repository');

class NotificationService {
  async listarNotificacoes(produtorId) {
    return await notificationRepository.listarPorProdutor(produtorId);
  }

  async marcarLida(notificacaoId, produtorId) {
    const result = await notificationRepository.marcarComoLida(notificacaoId, produtorId);
    
    if (!result) {
      const error = new Error('Notificação não encontrada ou não autorizada');
      error.statusCode = 404;
      throw error;
    }
    
    return { mensagem: 'Notificação marcada como lida com sucesso' };
  }
}

module.exports = new NotificationService();\end{lstlisting}
\vspace{0.5cm}
\subsubsection*{Arquivo: src/modules/order/order.controller.js}
\begin{lstlisting}[caption={src/modules/order/order.controller.js}]
// src/modules/order/order.controller.js
const orderService = require('./order.service');
const asyncHandler = require('../../middlewares/async.middleware');

class OrderController {
  listarPontosRetirada = asyncHandler(async (req, res) => {
    const pontos = await orderService.listarPontosRetirada();
    res.status(200).json(pontos);
  });

  fecharPedido = asyncHandler(async (req, res) => {
    const clienteId = req.user.id;
    const pedido = await orderService.fecharPedido(clienteId, req.body);
    res.status(201).json(pedido);
  });

  obterComprovante = asyncHandler(async (req, res) => {
    const { id } = req.params;
    const usuarioId = req.user.id;
    const perfil = req.user.perfil;
    const comprovante = await orderService.obterComprovante(id, usuarioId, perfil);
    res.status(200).json(comprovante);
  });

  listarHistorico = asyncHandler(async (req, res) => {
    const historico = await orderService.listarHistoricoCliente(req.user.id);
    res.status(200).json(historico);
  });

  listarPedidosProdutor = asyncHandler(async (req, res) => {
    const pedidos = await orderService.listarPedidosProdutor(req.user.id);
    res.status(200).json(pedidos);
  });

  atualizarStatus = asyncHandler(async (req, res) => {
    const { id } = req.params;
    const produtorId = req.user.id;
    const { status } = req.body;
    const resultado = await orderService.atualizarStatusPedido(id, produtorId, status);
    res.status(200).json(resultado);
  });
}

module.exports = new OrderController();\end{lstlisting}
\vspace{0.5cm}
\subsubsection*{Arquivo: src/modules/order/order.repository.js}
\begin{lstlisting}[caption={src/modules/order/order.repository.js}]
// src/modules/order/order.repository.js
const pool = require('../../config/db');

class OrderRepository {
  async listarPontosRetirada() {
    const query = `SELECT id, nome, endereco FROM pontos_retirada ORDER BY nome ASC`;
    const result = await pool.query(query);
    return result.rows;
  }

  async buscarEdicaoAtiva() {
    const query = `SELECT id, nome_identificador FROM edicoes_feira WHERE ativa = TRUE LIMIT 1`;
    const result = await pool.query(query);
    return result.rows[0] || null;
  }

  async buscarPontoRetiradaPorId(id) {
    const query = `SELECT id, nome, endereco FROM pontos_retirada WHERE id = $1`;
    const result = await pool.query(query, [id]);
    return result.rows[0] || null;
  }

  async buscarProdutosPorIds(ids) {
    const query = `
      SELECT id, produtor_id, nome, preco, ativo 
      FROM produtos 
      WHERE id = ANY($1::uuid[]) AND ativo = TRUE
    `;
    const result = await pool.query(query, [ids]);
    return result.rows;
  }

  async criarPedidoTransacional({ cliente_id, edicao_id, ponto_retirada_id, comprovante_codigo, status, valor_total, itens, notificacoes }) {
    const client = await pool.connect();
    try {
      await client.query('BEGIN');
      
      const insertPedidoQuery = `
        INSERT INTO pedidos (cliente_id, edicao_id, ponto_retirada_id, comprovante_codigo, status, valor_total)
        VALUES ($1, $2, $3, $4, $5, $6) 
        RETURNING id, comprovante_codigo, status, valor_total, criado_em;
      `;
      const pedidoRes = await client.query(insertPedidoQuery, [
        cliente_id, edicao_id, ponto_retirada_id, comprovante_codigo, status, valor_total
      ]);
      const pedido = pedidoRes.rows[0];

      const insertItemQuery = `
        INSERT INTO itens_pedido (pedido_id, produto_id, quantidade, preco_unitario)
        VALUES ($1, $2, $3, $4);
      `;
      for (const item of itens) {
        await client.query(insertItemQuery, [
          pedido.id, item.produto_id, item.quantidade, item.preco_unitario
        ]);
      }

      const insertNotificacaoQuery = `
        INSERT INTO notificacoes (produtor_id, pedido_id, mensagem)
        VALUES ($1, $2, $3);
      `;
      for (const notif of notificacoes) {
        await client.query(insertNotificacaoQuery, [
          notif.produtor_id, pedido.id, notif.mensagem
        ]);
      }

      await client.query('COMMIT');
      return pedido;
    } catch (error) {
      await client.query('ROLLBACK');
      throw error;
    } finally {
      client.release();
    }
  }

  async buscarComprovantePorId(pedidoId, usuarioId, perfil) {
    const pedidoQuery = `
      SELECT p.id, p.comprovante_codigo, p.status, p.valor_total, p.criado_em, p.cliente_id,
             pr.nome AS ponto_nome, pr.endereco AS ponto_endereco,
             ef.nome_identificador AS edicao_nome,
             u_cli.nome AS cliente_nome, u_cli.telefone AS cliente_telefone
      FROM pedidos p
      JOIN pontos_retirada pr ON p.ponto_retirada_id = pr.id
      JOIN edicoes_feira ef ON p.edicao_id = ef.id
      JOIN usuarios u_cli ON p.cliente_id = u_cli.id
      WHERE p.id = $1;
    `;
    const pedidoRes = await pool.query(pedidoQuery, [pedidoId]);
    const pedido = pedidoRes.rows[0];

    if (!pedido) return null;

    if (perfil === 'CLIENTE' && pedido.cliente_id !== usuarioId) {
      const error = new Error('Acesso não autorizado ao pedido');
      error.statusCode = 403;
      error.codigo = 'FORBIDDEN';
      throw error;
    }

    const itensQuery = `
      SELECT ip.quantidade, ip.preco_unitario, prod.nome AS produto_nome,
             u_prod.id AS produtor_id, u_prod.nome AS produtor_nome, u_prod.telefone AS produtor_telefone
      FROM itens_pedido ip
      JOIN produtos prod ON ip.produto_id = prod.id
      JOIN usuarios u_prod ON prod.produtor_id = u_prod.id
      WHERE ip.pedido_id = $1;
    `;
    const itensRes = await pool.query(itensQuery, [pedidoId]);
    pedido.itens = itensRes.rows;

    return pedido;
  }

  async listarPorCliente(clienteId) {
    const query = `
      SELECT p.id, p.comprovante_codigo, p.status, p.valor_total, p.criado_em,
             pr.nome AS ponto_nome, pr.endereco AS ponto_endereco
      FROM pedidos p
      JOIN pontos_retirada pr ON p.ponto_retirada_id = pr.id
      WHERE p.cliente_id = $1
      ORDER BY p.criado_em DESC;
    `;
    const result = await pool.query(query, [clienteId]);
    return result.rows;
  }

  async listarPorProdutor(produtorId) {
    const query = `
      SELECT DISTINCT p.id, p.comprovante_codigo, p.status, p.valor_total, p.criado_em,
             pr.nome AS ponto_nome,
             u_cli.nome AS cliente_nome, u_cli.telefone AS cliente_telefone
      FROM pedidos p
      JOIN pontos_retirada pr ON p.ponto_retirada_id = pr.id
      JOIN usuarios u_cli ON p.cliente_id = u_cli.id
      JOIN itens_pedido ip ON p.id = ip.pedido_id
      JOIN produtos prod ON ip.produto_id = prod.id
      WHERE prod.produtor_id = $1
      ORDER BY p.criado_em DESC;
    `;
    const result = await pool.query(query, [produtorId]);
    return result.rows;
  }

  async buscarItensDoProdutorNoPedido(pedidoId, produtorId) {
    const query = `
      SELECT ip.quantidade, ip.preco_unitario, prod.nome AS produto_nome
      FROM itens_pedido ip
      JOIN produtos prod ON ip.produto_id = prod.id
      WHERE ip.pedido_id = $1 AND prod.produtor_id = $2;
    `;
    const result = await pool.query(query, [pedidoId, produtorId]);
    return result.rows;
  }

  async verificarProdutorPertenceAoPedido(pedidoId, produtorId) {
    const query = `
      SELECT 1 FROM itens_pedido ip
      JOIN produtos prod ON ip.produto_id = prod.id
      WHERE ip.pedido_id = $1 AND prod.produtor_id = $2
      LIMIT 1;
    `;
    const result = await pool.query(query, [pedidoId, produtorId]);
    return result.rows.length > 0;
  }

  async buscarStatusAtual(pedidoId) {
    const query = `SELECT id, status, comprovante_codigo FROM pedidos WHERE id = $1;`;
    const result = await pool.query(query, [pedidoId]);
    return result.rows[0];
  }

  async atualizarStatus(pedidoId, novoStatus) {
    const query = `
      UPDATE pedidos 
      SET status = $1 
      WHERE id = $2 
      RETURNING id, status, comprovante_codigo;
    `;
    const result = await pool.query(query, [novoStatus, pedidoId]);
    return result.rows[0];
  }
}

module.exports = new OrderRepository();\end{lstlisting}
\vspace{0.5cm}
\subsubsection*{Arquivo: src/modules/order/order.routes.js}
\begin{lstlisting}[caption={src/modules/order/order.routes.js}]
// src/modules/order/order.routes.js
const express = require('express');
const orderController = require('./order.controller');
const { verifyToken, requireRole } = require('../../middlewares/auth.middleware');

const router = express.Router();

router.get('/pontos-retirada', orderController.listarPontosRetirada);
router.post('/pedidos', verifyToken, requireRole(['CLIENTE']), orderController.fecharPedido);
router.get('/pedidos/historico', verifyToken, requireRole(['CLIENTE']), orderController.listarHistorico);
router.get('/pedidos/produtor', verifyToken, requireRole(['PRODUTOR']), orderController.listarPedidosProdutor);
router.get('/pedidos/:id/comprovante', verifyToken, orderController.obterComprovante);
router.patch('/pedidos/:id/status', verifyToken, requireRole(['PRODUTOR']), orderController.atualizarStatus);

module.exports = router;\end{lstlisting}
\vspace{0.5cm}
\subsubsection*{Arquivo: src/modules/order/order.service.js}
\begin{lstlisting}[caption={src/modules/order/order.service.js}]
// src/modules/order/order.service.js
const orderRepository = require('./order.repository');

class OrderService {
  async listarPontosRetirada() {
    return await orderRepository.listarPontosRetirada();
  }

  async fecharPedido(clienteId, { ponto_retirada_id, itens }) {
    if (!ponto_retirada_id) {
      const error = new Error('Ponto de retirada é obrigatório');
      error.statusCode = 400;
      throw error;
    }

    if (!Array.isArray(itens) || itens.length === 0) {
      const error = new Error('O carrinho não pode estar vazio');
      error.statusCode = 400;
      throw error;
    }

    const ponto = await orderRepository.buscarPontoRetiradaPorId(ponto_retirada_id);
    if (!ponto) {
      const error = new Error('Ponto de retirada inválido');
      error.statusCode = 404;
      throw error;
    }

    const edicao = await orderRepository.buscarEdicaoAtiva();
    if (!edicao) {
      const error = new Error('Não há nenhuma edição de feira ativa no momento');
      error.statusCode = 400;
      throw error;
    }

    const idsProdutos = itens.map(i => i.produto_id);
    const dbProdutos = await orderRepository.buscarProdutosPorIds(idsProdutos);

    if (dbProdutos.length !== idsProdutos.length) {
      const error = new Error('Um ou mais produtos não estão mais disponíveis');
      error.statusCode = 400;
      throw error;
    }

    let valor_total = 0;
    const itensMapeados = [];
    const produtoresSet = new Set();

    for (const reqItem of itens) {
      if (reqItem.quantidade <= 0) {
        const error = new Error('Quantidade inválida para o produto ' + reqItem.produto_id);
        error.statusCode = 400;
        throw error;
      }
      
      const dbProd = dbProdutos.find(p => p.id === reqItem.produto_id);
      const precoUnitario = Number(dbProd.preco);
      valor_total += precoUnitario * reqItem.quantidade;
      
      itensMapeados.push({
        produto_id: reqItem.produto_id,
        quantidade: reqItem.quantidade,
        preco_unitario: precoUnitario
      });

      produtoresSet.add(dbProd.produtor_id);
    }

    const comprovante_codigo = 'SISF-' + Math.floor(1000 + Math.random() * 9000);
    const status = 'RECEBIDO';

    const notificacoes = Array.from(produtoresSet).map(produtor_id => ({
      produtor_id,
      mensagem: 'Você tem um novo pedido com itens a serem preparados'
    }));

    const pedido = await orderRepository.criarPedidoTransacional({
      cliente_id: clienteId,
      edicao_id: edicao.id,
      ponto_retirada_id,
      comprovante_codigo,
      status,
      valor_total,
      itens: itensMapeados,
      notificacoes
    });

    return pedido;
  }

  async obterComprovante(pedidoId, usuarioId, perfil) {
    const comprovante = await orderRepository.buscarComprovantePorId(pedidoId, usuarioId, perfil);
    if (!comprovante) {
      const error = new Error('Pedido não encontrado');
      error.statusCode = 404;
      throw error;
    }
    
    comprovante.valor_total = Number(comprovante.valor_total);
    comprovante.itens = comprovante.itens.map(item => ({
      ...item,
      preco_unitario: Number(item.preco_unitario)
    }));
    
    return comprovante;
  }

  async listarHistoricoCliente(clienteId) {
    const historico = await orderRepository.listarPorCliente(clienteId);
    return historico.map(p => ({
      ...p,
      valor_total: Number(p.valor_total)
    }));
  }

  async listarPedidosProdutor(produtorId) {
    const pedidos = await orderRepository.listarPorProdutor(produtorId);
    
    for (const pedido of pedidos) {
      pedido.valor_total = Number(pedido.valor_total);
      pedido.itens = await orderRepository.buscarItensDoProdutorNoPedido(pedido.id, produtorId);
    }
    
    return pedidos;
  }

  async atualizarStatusPedido(pedidoId, produtorId, novoStatus) {
    const statusValidos = ['RECEBIDO', 'EM_PREPARACAO', 'ENTREGUE'];
    if (!statusValidos.includes(novoStatus)) {
      const error = new Error('Status inválido');
      error.statusCode = 400;
      throw error;
    }

    const pedido = await orderRepository.buscarStatusAtual(pedidoId);
    if (!pedido) {
      const error = new Error('Pedido não encontrado');
      error.statusCode = 404;
      throw error;
    }

    const pertence = await orderRepository.verificarProdutorPertenceAoPedido(pedidoId, produtorId);
    if (!pertence) {
      const error = new Error('Acesso não autorizado para alterar este pedido');
      error.statusCode = 403;
      throw error;
    }

    const statusAtual = pedido.status;
    if (statusAtual === 'RECEBIDO' && novoStatus !== 'EM_PREPARACAO') {
      const error = new Error('Transição inválida: pedidos recebidos só podem avançar para EM_PREPARACAO');
      error.statusCode = 400;
      throw error;
    }
    if (statusAtual === 'EM_PREPARACAO' && novoStatus !== 'ENTREGUE') {
      const error = new Error('Transição inválida: pedidos em preparação só podem avançar para ENTREGUE');
      error.statusCode = 400;
      throw error;
    }
    if (statusAtual === 'ENTREGUE') {
      const error = new Error('Transição inválida: pedido já se encontra finalizado como ENTREGUE');
      error.statusCode = 400;
      throw error;
    }

    const atualizado = await orderRepository.atualizarStatus(pedidoId, novoStatus);
    
    return { 
      sucesso: true, 
      mensagem: `Status do pedido alterado para ${novoStatus} com sucesso.`,
      status: atualizado.status
    };
  }
}

module.exports = new OrderService();\end{lstlisting}
\vspace{0.5cm}
\subsubsection*{Arquivo: src/modules/report/report.controller.js}
\begin{lstlisting}[caption={src/modules/report/report.controller.js}]
// src/modules/report/report.controller.js
const reportService = require('./report.service');
const asyncHandler = require('../../middlewares/async.middleware');

class ReportController {
  obterRelatorioVendas = asyncHandler(async (req, res) => {
    const relatorio = await reportService.gerarRelatorioVendas(req.user.id);
    res.status(200).json(relatorio);
  });
}

module.exports = new ReportController();\end{lstlisting}
\vspace{0.5cm}
\subsubsection*{Arquivo: src/modules/report/report.repository.js}
\begin{lstlisting}[caption={src/modules/report/report.repository.js}]
// src/modules/report/report.repository.js
const pool = require('../../config/db');

class ReportRepository {
  async obterMetricasGerais(produtorId) {
    const query = `
      SELECT 
        COALESCE(SUM(ip.quantidade * ip.preco_unitario), 0)::numeric(10,2) AS total_faturado,
        COUNT(DISTINCT ip.pedido_id)::int AS total_pedidos
      FROM itens_pedido ip
      JOIN produtos p ON ip.produto_id = p.id
      WHERE p.produtor_id = $1;
    `;
    const result = await pool.query(query, [produtorId]);
    return result.rows[0];
  }

  async obterItensMaisVendidos(produtorId) {
    const query = `
      SELECT 
        p.nome AS produto_nome,
        SUM(ip.quantidade)::int AS quantidade_vendida,
        SUM(ip.quantidade * ip.preco_unitario)::numeric(10,2) AS subtotal
      FROM itens_pedido ip
      JOIN produtos p ON ip.produto_id = p.id
      WHERE p.produtor_id = $1
      GROUP BY p.nome
      ORDER BY quantidade_vendida DESC;
    `;
    const result = await pool.query(query, [produtorId]);
    return result.rows;
  }
}

module.exports = new ReportRepository();\end{lstlisting}
\vspace{0.5cm}
\subsubsection*{Arquivo: src/modules/report/report.routes.js}
\begin{lstlisting}[caption={src/modules/report/report.routes.js}]
// src/modules/report/report.routes.js
const express = require('express');
const reportController = require('./report.controller');
const { verifyToken, requireRole } = require('../../middlewares/auth.middleware');

const router = express.Router();

router.get('/vendas/produtor', verifyToken, requireRole(['PRODUTOR']), reportController.obterRelatorioVendas);

module.exports = router;\end{lstlisting}
\vspace{0.5cm}
\subsubsection*{Arquivo: src/modules/report/report.service.js}
\begin{lstlisting}[caption={src/modules/report/report.service.js}]
// src/modules/report/report.service.js
const reportRepository = require('./report.repository');

class ReportService {
  async gerarRelatorioVendas(produtorId) {
    const [metricas, itens] = await Promise.all([
      reportRepository.obterMetricasGerais(produtorId),
      reportRepository.obterItensMaisVendidos(produtorId)
    ]);

    return {
      total_faturado: parseFloat(metricas.total_faturado) || 0,
      total_pedidos: parseInt(metricas.total_pedidos, 10) || 0,
      itens_mais_vendidos: itens.map(item => ({
        produto_nome: item.produto_nome,
        quantidade_vendida: parseInt(item.quantidade_vendida, 10),
        subtotal: parseFloat(item.subtotal)
      }))
    };
  }
}

module.exports = new ReportService();\end{lstlisting}
\vspace{0.5cm}
